\section{Methodik}\label{methodik}

Der hier beschriebene Versuchsaufbau setzt die zuvor identifizierte Forschungsluecke um.
Wichtig war dabei, jeden Schritt von der Datenaufbereitung bis zur \shap{}-Analyse so
festzuhalten, dass sich die Ergebnisse exakt wiederholen lassen, denn ohne diese
Nachvollziehbarkeit waere eine erklaerbare Angriffserkennung wenig wert.

\subsection{Datenbasis und Vorverarbeitung}

Als Datenbasis dient der CICIoT2023-Datensatz~\cite{neto2023ciciot2023}, der Netzwerkverkehr
von 105 realen \iot{}-Geraeten unter 33 verschiedenen Angriffsbedingungen enthaelt.
Die Daten liegen als vorverarbeitete CSV-Dateien (Comma-Separated Values) mit 39 extrahierten
Netzwerkmerkmalen vor sowie einer Klassenspalte (Label).

\textbf{Hinweis zur Feature-Anzahl.}
Raturi et al.\ \cite{raturi2026exploratory} und Alharby~\cite{alharby2025ciciot} benennen 46
Features fuer den CICIoT2023-Datensatz.
Die vorliegenden CSV-Dateien enthalten jedoch 39 Merkmale.
Diese Abweichung wird hier dokumentiert und unveraendert uebernommen.
\pca{} in Pipeline~A reduziert demnach von 39 auf 16 Hauptkomponenten.

Die 33 Angriffsklassen werden zu 8 uebergeordneten Kategorien zusammengefasst:
DDoS, DoS, Mirai, Reconnaissance, Spoofing, Brute-Force, Web-based Attacks und Benign.
Instanzen der Klasse BACKDOOR\_MALWARE, die in keiner dieser 8 Kategorien enthalten ist
und unter 0{,}01~Prozent des Datenvolumens ausmacht, werden vor dem Training entfernt.
Das bereinigte Dataset umfasst 838.602 Instanzen, davon 58{,}4~Prozent DoS/DDoS-Verkehr.

\textbf{Stratifiziertes Sampling.}
Da das vollstaendige Dataset von ca.\ 8{,}5~GB die verfuegbare Arbeitsspeicherkapazitaet
ueberschreitet, wird ein stratifiziertes Sampling-Verfahren eingesetzt.
Aus jeder der 63 CSV-Dateien werden je 500 Instanzen pro Klasse entnommen, sodass alle
Dateien gleichmaessig zum Trainingsdatensatz beitragen.
Jede Datei wird dabei vollstaendig eingelesen, bevor die Stichprobe gezogen und die Datei
wieder aus dem Speicher entfernt wird; so wird keine Zeile von vornherein ausgeschlossen, und
der Speicherbedarf bleibt beschraenkt.
Dadurch bleiben alle Klassen im Trainingsdatensatz vertreten, und das ausgepraegte
Klassenungleichgewicht des Gesamtdatensatzes wird abgemildert.

Unendliche Werte (Infinity), die im Rohdatensatz auftreten, werden vor dem Preprocessing
durch fehlende Werte (NaN) ersetzt.
Ein anschliessender Median-Imputer fuellt fehlende Werte mit dem Spaltenmedian auf.
Diese Berechnung erfolgt ausschliesslich auf den Trainingsdaten, um eine Datenleckage zu
vermeiden.

\textbf{Aufteilung.}
Der Datensatz wird in 80~Prozent Trainings- und 20~Prozent Testdaten aufgeteilt, stratifiziert
nach Kategorie mit festem Random-Seed 42 fuer vollstaendige Reproduzierbarkeit.

\subsection{Zwei parallele Preprocessing-Pipelines}

Zwei Preprocessing-Pipelines werden verglichen:

\textbf{Pipeline~A} setzt eine \pca{}-Dimensionsreduktion ein.
Sie besteht aus drei Schritten: SimpleImputer (Medianstrategie), StandardScaler und
\pca{} (39 auf 16 Hauptkomponenten).
Dieser Ansatz repliziert das Vorgehen von Raturi et al.\ \cite{raturi2026exploratory}, die
\pca{} als Dimensionsreduktionstechnik einsetzen, und von Alharby~\cite{alharby2025ciciot},
der \pca{} auf demselben Datensatz anwendet.

\textbf{Pipeline~B} verzichtet vollstaendig auf \pca{}.
Sie besteht aus SimpleImputer und StandardScaler.
Alle 39 Originalmerkmale bleiben erhalten.
Da \shap{}-Werte auf den originalen Merkmalsnamen basieren, sind die Erklaerungen fuer
Sicherheitsanalysten unmittelbar semantisch interpretierbar~\cite{mohale2025xai}.

\subsection{Klassifikationsarchitektur}

Als Klassifikationsarchitektur wird der zweistufige hierarchische Ansatz aus Raturi et al.\
\cite{raturi2026exploratory} uebernommen.
Abbildung~\ref{fig:architektur} fasst die gesamte Verarbeitungskette von den Rohdaten bis zur
\shap{}-Analyse zusammen.
In Stufe~1 erfolgt eine binaere Klassifikation zwischen DoS/DDoS-Verkehr und
Nicht-DoS-Verkehr.
In Stufe~2 werden die Nicht-DoS-Instanzen in die verbleibenden Kategorien eingeordnet:
Mirai, Reconnaissance, Spoofing, Brute-Force, Web-based Attacks und Benign.
Diese Trennung adressiert das Klassenungleichgewicht strukturell, da die dominanten DoS- und
DDoS-Klassen in einer separaten binaeren Stufe behandelt werden, bevor die
Mehrklassenklassifikation der zahlenmaessig unterrepraesentierten Kategorien erfolgt.

Stufe~2 wird ausschliesslich auf den wahren Nicht-DoS-Trainingsinstanzen trainiert.
Dadurch wird verhindert, dass Klassifikationsfehler aus Stufe~1 in das Training von Stufe~2
einfliessen.
Zur Vorhersagezeit reicht Stufe~1 nur die als Nicht-DoS eingestuften Instanzen an Stufe~2
weiter, waehrend alle uebrigen direkt als DoS/DDoS ausgegeben werden.

\begin{figure}[h]
\centering
\begin{tikzpicture}[
  box/.style={draw, rounded corners, align=center, inner sep=3pt,
              font=\footnotesize, text width=0.78\columnwidth, minimum height=6mm},
  arr/.style={-{Latex[length=2mm]}, thick}
]
\node[box] (data) {CICIoT2023 (CSV)\\ 39 Merkmale, 838.602 Instanzen};
\node[box, below=4mm of data] (prep) {Vorverarbeitung\\ Median-Imputer, StandardScaler\\
  (Pipeline~A zusaetzlich \pca{}: 39 $\to$ 16)};
\node[box, below=4mm of prep] (s1) {Stufe~1 (binaer, \gb{})\\ DoS/DDoS vs.\ Nicht-DoS};
\node[box, below=4mm of s1] (s2) {Stufe~2 (Mehrklasse, \gb{})\\ Mirai, Reconnaissance, Spoofing,
  Brute-Force, Web-based, Benign};
\node[box, below=4mm of s2] (shap) {SHAP\\ TreeExplainer (Stufe~1),
  PermutationExplainer (Stufe~2)};
\draw[arr] (data) -- (prep);
\draw[arr] (prep) -- (s1);
\draw[arr] (s1) -- (s2) node[midway, right, font=\scriptsize] {nur Nicht-DoS};
\draw[arr] (s2) -- (shap);
\end{tikzpicture}
\caption{Verarbeitungskette der Arbeit. Beide Pipelines teilen dieselbe Struktur; Pipeline~A
schaltet zusaetzlich eine \pca{}-Reduktion vor. Stufe~2 verarbeitet ausschliesslich
Nicht-DoS-Instanzen.}
\label{fig:architektur}
\end{figure}

\subsection{Modellwahl}

Als Hauptmodell wird \gb{} implementiert.
Erstens erzielt \gb{} in Raturi et al.\ \cite{raturi2026exploratory} die hoechste \ba{} in
der Mehrklassenklassifikation auf dem CICIoT2023-Datensatz.
Zweitens sind \gb{}-Modelle vollstaendig kompatibel mit dem SHAP-TreeExplainer, der native
und exakte \shap{}-Werte ohne Approximation berechnet~\cite{mohale2025xai,hermosilla2025xai_forensic}.

\subsection{Evaluationsmetriken}

Die primaere Bewertungsmetrik ist die \ba{}, die in Raturi et al.\ \cite{raturi2026exploratory}
als geeignete Metrik fuer den unausgewogenen CICIoT2023-Datensatz vorgeschlagen wird.
Sie ist definiert als arithmetisches Mittel der klassenspezifischen Recall-Werte und ist robust
gegenueber Klassenungleichgewicht.
Hosseini et al.\ \cite{hosseini2025imbalance} begruenden empirisch, warum einfache Accuracy
bei unausgewogenen \ids{}-Datensaetzen als alleinige Metrik ungenuegend ist.
Als sekundaere Metriken werden der F1-Score makro-gemittelt, Precision und Recall pro Klasse
sowie Konfusionsmatrizen erhoben.

\subsection{Ablation Study}

Die Ablation Study vergleicht Pipeline~A (mit \pca{}) und Pipeline~B (ohne \pca{}) unter
identischen Modellparametern.
Untersucht werden der Einfluss der \pca{} auf die \ba{} beider Klassifikationsstufen sowie
auf die \shap{}-Feature-Importance-Ranglisten.
Alharby~\cite{alharby2025ciciot} liefert Vergleichswerte fuer \pca{}-basierte Klassifikation
auf dem CICIoT2023-Datensatz.

Da \pca{}-Hauptkomponenten lineare Kombinationen der Originalmerkmale sind, verliert jede
Komponente die direkte semantische Bedeutung eines konkreten Netzwerkmerkmals.
Die Ablation Study prueft, ob dieser Verlust an Interpretierbarkeit durch eine verbesserte
Modellleistung aufgewogen wird.

\subsection{SHAP-Integration}

\shap{} wird ausschliesslich als post-hoc-Erklaerungskomponente eingesetzt.
Die Analyse erfolgt auf zwei Ebenen in Anlehnung an Mohale und Obagbuwa~\cite{mohale2025xai}.

\textbf{Explainer-Wahl.}
Fuer Stufe~1 (binaere Klassifikation) wird der \texttt{TreeExplainer} eingesetzt, der
exakte Shapley-Werte fuer baumbasierte Modelle ohne Approximation
berechnet~\cite{mohale2025xai}.
Fuer Stufe~2 (Mehrklassenklassifikation mit mehr als zwei Klassen) ist der
\texttt{TreeExplainer} in \shap{}~0.52 fuer \texttt{GradientBoostingClassifier} nicht
verfuegbar.
Daher kommt der \texttt{PermutationExplainer} mit \texttt{predict\_proba} als
Modellschnittstelle zum Einsatz, der modellunabhaengige Shapley-Werte berechnet und die
wissenschaftliche Gueltigkeit der Erklaerungen
erhaelt~\cite{mohale2025xai,hermosilla2025xai_forensic}.

Auf \textbf{globaler Ebene} werden Feature-Importance-Ranglisten pro Angriffsklasse erzeugt,
basierend auf dem Betragsmittelwert der \shap{}-Werte ueber alle Testinstanzen
($\overline{|\phi_i|}$).

Auf \textbf{lokaler Ebene} werden Einzelfallanalysen ausgewaehlter Fehlklassifikationen
durchgefuehrt (Waterfall-Plots).

Die Qualitaet der \shap{}-Erklaerungen wird entlang der drei von Hermosilla et al.\
\cite{hermosilla2025xai_forensic} benannten Dimensionen quantitativ bewertet, hier wie folgt
operationalisiert.
Die Faithfulness (Fidelitaet) ist die Pearson-Korrelation zwischen den \shap{}-Werten der
vorhergesagten Klasse und der tatsaechlichen Aenderung der Klassenwahrscheinlichkeit, wenn ein
Merkmal auf seinen Hintergrund-Median gesetzt wird.
Die Stabilitaet (Konsistenz) ist die mittlere Kosinus-Aehnlichkeit zwischen der Erklaerung einer
Instanz und den Erklaerungen kleiner gaussscher Stoerungen; die Jaccard-Aehnlichkeit misst die
Rangstabilitaet der fuenf wichtigsten Merkmale unter denselben Stoerungen.
Alle drei Masse werden deterministisch auf den Stufe-2-Erklaerungen beider Pipelines berechnet.
Die Comprehensibility wird ergaenzend ueber die semantische Validierung adressiert, die die
Feature-Importance-Ranglisten mit den dokumentierten Netzwerksignaturen der Angriffstypen
abgleicht~\cite{neto2023ciciot2023}.
